\subsection{Symmetric rational fractions}

Let K be a commutative field and \(X_{1}, X_{2}, \ldots, X_{n}\) indeterminates. For every \(\sigma \in S_{n}\) we have defined in No.~1 (IV, p.~61) an automorphism \(\varphi_{\sigma}\) of \(K[X_{1}, X_{2}, \ldots, X_{n}]\) . This automorphism has a unique extension \(\psi_{\sigma}\) to the field \(K(X_{1}, \ldots, X_{n})\) , and \(\sigma \mapsto \psi_{\sigma}\) is an injective homomorphism of \(S_{n}\) into the automorphism group of \(K(X_{1}, \ldots, X_{n})\) . For each \(f \in K(X_{1}, \ldots, X_{n})\) we have \((\psi_{\sigma} f) (X_{1}, \ldots, X_{n}) = f(X_{\sigma(1)}, \ldots, X_{\sigma(n)})\) . The rational fractions f such that \(\psi_{\sigma}(f) = f\) for all \(\sigma \in S_{n}\) are called symmetric rational fractions. The set of all symmetric rational fractions in \(X_{1}, \ldots, X_{n}\) is a subfield of \(K(X_{1}, \ldots, X_{n})\) .

PROPOSITION 2. --- The field of symmetric rational fractions in \(X_{1}, \ldots, X_{n}\) is the field of fractions of the ring of symmetric polynomials in \(X_{1}, \ldots, X_{n}\) .

Let \(f \in \mathbf{K}(\mathbf{X}_1, \ldots, \mathbf{X}_n)\) be a symmetric rational fraction, and let \(u_1, v_1\) be two elements of \(\mathbf{K}[\mathbf{X}_1, \ldots, \mathbf{X}_n]\) such that \(f = \frac{u_1}{v_1}\) . We put \(v = \prod_{\sigma \in \mathfrak{S}_n} \psi_\sigma(v_1) \in \mathbf{K}[\mathbf{X}_1, \ldots, \mathbf{X}_n]\) and \(u = fv \in \mathbf{K}[\mathbf{X}_1, \ldots, \mathbf{X}_n]\) ; then \(v\) is symmetric, hence \(u\) is symmetric, because \(f\) is, and \(f = \frac{u}{v}\) , whence the result.

COROLLARY. --- Let \(s_{1}, s_{2}, \ldots, s_{n}\) be the elementary symmetric polynomials in \(X_{1}, \ldots, X_{n}\) . For every rational fraction \(g \in \mathbf{K}(S_{1}, S_{2}, \ldots, S_{n})\) the sequence \((s_{1}, s_{2}, \ldots, s_{n})\) is substitutable in g and the mapping \(g \mapsto g(s_{1}, s_{2}, \ldots, s_{n})\) is an isomorphism of \(\mathbf{K}(S_{1}, S_{2}, \ldots, S_{n})\) onto the field of symmetric rational fractions in \(X_{1}, \ldots, X_{n}\) .

This follows from Prop. 2 and Th. 1 of IV, p.~62.

